\documentclass[11pt]{article}
\usepackage{amsmath, amssymb, geometry, hyperref}
\geometry{margin=1in}

\title{\textbf{Cosmological Relaxation: A Mechanical Interpretation of the Early Universe Phase Transition}}
\author{Research Manuscript Supplement}
\date{2026}

\begin{document}
\maketitle

\begin{abstract}
We propose a model wherein the early universe is described as a transition from a state of absolute metric rigidity to a viscoelastic elastic regime. By applying the Tensor Field saturation factor $\xi$, we suggest that the Big Bang represents a "Universal Relaxation Event" occurring as the total mass-energy density $\rho_U$ dropped below the vacuum's structural modulus.
\end{abstract}

\section{The Frozen Core and Universal Relaxation}
In this framework, the initial state of the universe is modeled as a "Frozen Core" where $\xi \to 0$. In this regime, the mass-energy density $\rho_U$ creates an absolute saturation of the field:
\begin{equation}
\xi_{universal} = \frac{8\pi u_\nu}{8\pi u_\nu + \rho_U c^2}
\end{equation}
This state implies a metric that is effectively incompressible, where constituent displacement and wave propagation ($c$) are inhibited by the total saturation of the medium.

\section{The Shock-Front Hypothesis}
As expansion initiated a decrease in $\rho_U$, the field transitioned from a rigid state to an elastic one. We suggest that this transition characterizes a non-linear mechanical event, analogous to a shockwave in a fluid medium. We propose that the Cosmic Microwave Background (CMB) may be re-evaluated as the thermalized relic of this initial "Sonic Boom" or shock front, marking the moment the vacuum regained its conductive elasticity.

\section{Disclosure and Sentiment}
The author posits that this model offers a mechanical continuity between the relativistic limits of particles and the macroscopic evolution of the cosmos. This framework is intended to supplement, rather than replace, current inflationary models by providing a constitutive material basis for the expansion.

\begin{thebibliography}{9}
\bibitem{Guth1981} Guth, A. H. (1981). Inflationary universe: A possible solution to the horizon and flatness problems. \textit{Phys. Rev. D}.
\bibitem{Penrose2010} Penrose, R. (2010). \textit{Cycles of Time: An Extraordinary New View of the Universe}.
\end{thebibliography}

\end{document}